\chapter{Plano e Especificação de Testes de Software}
\label{chap:plano_testes}

\section{Estratégia e Abordagem de Testes}
A garantia da qualidade do software segue as diretrizes da norma internacional \textbf{ISO/IEC/IEEE 29119}. A estratégia apoia-se no modelo de Pirâmide de Testes, priorizando uma ampla base de testes unitários automatizados, complementada por testes de integração de API e testes ponta a ponta (E2E) para fluxos críticos.

\begin{itemize}
    \item \textbf{Testes Unitários:} Validação isolada de regras de negócio, entidades de domínio e métodos de serviço;
    \item \textbf{Testes de Integração:} Validação dos endpoints HTTP, transações com o banco de dados PostgreSQL e geração de tokens JWT;
    \item \textbf{Testes de Sistema / E2E:} Simulação de fluxos completos do usuário a partir da interface React;
    \item \textbf{Testes de Desempenho e Carga:} Medição de latência e vazão sob requisições concorrentes via k6 ou JMeter.
\end{itemize}

% PLACEHOLDER STARUML v7.0: Pirâmide e Estratégia de Testes
% Ferramenta: StarUML v7.0 -> Model -> Add Diagram -> Component Diagram ou Package Diagram
% Arquivo: Imagens/test_piramide_estrategia.png (Exportar em PNG 300 DPI ou PDF vetorial)
% Descomente a linha abaixo apos salvar o arquivo no diretorio Imagens/:
% \incluirdiagrama{Imagens/test_piramide_estrategia.png}{Pirâmide e Níveis de Testes Automatizados}{fig:test_piramide}
\begin{figure}[htbp]
    \centering
    \fbox{\parbox{0.9\textwidth}{\centering\vspace{1cm}%
        \textbf{[Pirâmide e Níveis de Teste -- StarUML v7.0]}\\%
        \textit{Modelar no StarUML v7.0 e exportar para: \texttt{Imagens/test\_piramide\_estrategia.png}}\\%
        \small Menu: \texttt{File -> Export Diagram as -> PNG...} (Fundo branco, alta resolução 300 DPI)\\%
        \vspace{0.3cm}%
        \footnotesize Para ativar a imagem real, descomente no arquivo \texttt{.tex}:\\
        \texttt{\textbackslash incluirdiagrama\{Imagens/test\_piramide\_estrategia.png\}\{Pirâmide de Testes\}\{fig:test\_piramide\}}%
    \vspace{1cm}}}
    \caption{Pirâmide de Testes e Estratégia de Cobertura (Placeholder StarUML v7.0)}
    \label{fig:test_piramide}
\end{figure}

\section{Ambientes e Ferramentas de Teste}
\begin{table}[htbp]
\caption{Ferramentas do Ecossistema de Testes}
\label{tab:ferramentas_teste}
\centering
\small
\begin{tabularx}{\textwidth}{@{} l l Y @{}}
\toprule
\textbf{Nível de Teste} & \textbf{Ferramenta} & \textbf{Escopo e Utilização} \\
\midrule
Unitário (Back-end) & Pytest + pytest-mock & Teste de regras de negócio e validações Pydantic. \\
Integração (API) & HTTPX + pytest-asyncio & Testes assíncronos das rotas FastAPI em banco de teste. \\
Unitário (Front-end) & Vitest + React Testing Library & Renderização de componentes e simulação de cliques. \\
Carga / Estresse & k6 & Teste de pico de acessos simultâneos ao agendamento. \\
\bottomrule
\end{tabularx}
\end{table}

\section{Critérios de Entrada, Suspensão e Aceite}
\begin{itemize}
    \item \textbf{Critério de Aceite de Versão:} Cobertura mínima de código $\ge 80\%$ das linhas de negócio, $100\%$ dos testes de regressão passando no pipeline CI e zero bugs de severidade Crítica ou Alta abertos;
    \item \textbf{Critério de Suspensão:} Falha que impeça o build da imagem Docker ou indisponibilidade do banco de dados nos testes automatizados.
\end{itemize}

\section{Catálogo de Casos de Teste (CT)}
Os casos de teste foram especificados de forma determinística e reproduzível, cobrindo cenários nominais e de exceção. A Tabela~\ref{tab:caso_teste_ct01} exemplifica o caso de teste referente ao agendamento de consultas.

\begin{table}[H]
\caption{Caso de Teste CT01 -- Validação de Conflito de Horário}
\label{tab:caso_teste_ct01}
\centering
\small
\renewcommand{\arraystretch}{1.25}
\begin{tabularx}{\textwidth}{@{} l Y @{}}
\toprule
\multicolumn{2}{@{}l}{\textbf{CT01 -- Tentativa de Agendamento em Horário Ocupado}} \\
\midrule
\textbf{Requisito / UC} & RF02 / UC01 \\
\textbf{Pré-condições} & Profissional P1 possui agendamento confirmado para o dia D às 14:00. \\
\textbf{Procedimento} & 
1. Realizar requisição POST para \texttt{/api/v1/agendamentos};\newline
2. Enviar payload com \texttt{profissional\_id = P1} e \texttt{data\_hora = D 14:00};\newline
3. Capturar o código de status HTTP e o corpo da resposta. \\
\textbf{Resultado Esperado} & Status HTTP 409 (Conflict) com mensagem de erro informando indisponibilidade do horário. Nenhuma linha é duplicada no banco. \\
\textbf{Critério de Sucesso} & Passagem automática no teste de integração da API. \\
\bottomrule
\end{tabularx}
\end{table}

\section{Matriz de Rastreabilidade (Requisitos vs Testes)}
A Tabela~\ref{tab:rastreabilidade_testes} consolida o mapeamento bidirecional entre os requisitos e as suítes de teste automatizadas, atestando a cobertura completa de validação.

\begin{table}[H]
\caption{Matriz de Rastreabilidade entre Requisitos e Casos de Teste}
\label{tab:rastreabilidade_testes}
\centering
\small
\renewcommand{\arraystretch}{1.25}
\begin{tabularx}{\textwidth}{@{} l l Y c @{}}
\toprule
\textbf{Requisito} & \textbf{Descrição Sintética} & \textbf{Casos de Teste Vinculados} & \textbf{Resultado} \\
\midrule
RF01 & Autenticação e JWT & CT-AUTH-01, CT-AUTH-02, CT-AUTH-03 & \badgeconcluido \\
RF02 & Agendamento de Horários & CT01, CT02, CT-LOAD-01 & \badgeconcluido \\
RF03 & Prontuário Imutável & CT-PEP-01, CT-PEP-02 & \badgeemprogresso \\
RNF01 & Tempo de Resposta $\le 2{,}0$s & CT-LOAD-01 (k6 estresse) & \badgeconcluido \\
\bottomrule
\end{tabularx}
\end{table}
